% Integración sucesora 2026-09-26: el lector longitudinal precede a G.
% La realización circular anterior se conserva como corolario, no como generador.
\section{Acción, longitud estructural, gravedad e información térmica}
\label{v:ant:sec:gravity}

\subsection{Retorno registrado y construcción de la longitud}

Escribimos $\alpha_{\HMT}$ para la coordenada dodecafásica ya generada
---denotada $\alpha$ en las secciones anteriores---. No interviene aquí
ningún valor metrológico de $G$ ni una longitud de Planck tabulada. El
registro marcado reúne
\begin{equation}
 N=12\cdot4\cdot120=5760
 \label{v:eq:cardinal-longitud-alpha}
\end{equation}
posiciones. Sean $u_\Omega=N^{-1/2}\mathbf1_\Omega$,
$Q_\Omega=I-u_\Omega u_\Omega^*$ y
$\mathsf B:V_U\to V_\Omega$ el retorno de Hadamard tipado, con
\begin{equation}
 \mathsf B^*\mathsf B=\mathsf B\mathsf B^*=\frac14I,
 \qquad \mathsf C=Q_\Omega\mathsf B=\mathsf BQ_U.
 \label{v:eq:retorno-hadamard-longitud}
\end{equation}
Para el pinzamiento diagonal
$\Delta_\Omega(A)=\sum_{i\in\Omega}E_iAE_i$, se obtiene
\begin{equation}
 \Delta_\Omega(\mathsf C\mathsf C^*)
 =\frac{N-1}{4N}I=r_NI,
 \qquad
 r_N=\frac{5759}{23040},
 \label{v:eq:residuo-longitud-alpha}
\end{equation}
mientras el canal complementario conserva $1/(4N)$. En efecto,
$\mathsf C\mathsf C^*=Q_\Omega/4$ y cada entrada diagonal de
$Q_\Omega$ vale $1-1/N$. El pinzamiento determina la respuesta
longitudinal; la inscripción ampliada y su complemento conservan la
memoria que la lectura diagonal no distingue.

Sea $\mathcal L_L$ la recta orientada de longitud y sea $L_*>0$ su
sección de referencia transportada por la cocadena longitudinal. Esta
sección no se identifica silenciosamente con la base nativa de arista:
su relación con el reloj se fija después por la realización circular.
La norma del conúcleo $H_4$ ya construida aporta
$\mathrm N_{H_4}(\alpha_{\HMT})=\alpha_{\HMT}^{16}$. La composición
lineal del retorno registrado determina, por tanto,
\begin{equation}
 \boxed{
 L_\alpha=\alpha_{\HMT}^{16}r_NL_*
 =\alpha_{\HMT}^{16}\frac{5759}{23040}L_* .}
 \label{v:eq:longitud-alpha}
\end{equation}
La potencia, el coeficiente de incidencia y la sección dimensional tienen
orígenes separados y explícitos. $L_\alpha$ es una salida longitudinal
anterior a la lectura gravitatoria.

\subsection{Composición operatoria y determinación única de gravedad}

Sea $\mathcal H\ne\{0\}$ una fibra de Hilbert finita y sea $Y=UY_0U^*>0$
el carácter positivo e invertible de una ruta, con su transporte, hoja,
orientación y memoria conservados. Para una sección de acción
$\hbar_a>0$, $a\in\{\mathrm{pre},\mathrm{ret}\}$, y la velocidad
$c=c_{\mathrm{int}}=c_{\mathrm{vac}}=\widehat c\,u_C>0$ ya construida en
\eqref{v:eq:identidad-velocidad}, pongamos
\begin{equation}
 Q_{L,a}=\frac{\hbar_ac}{L_\alpha},\qquad
 H_a=Q_{L,a}Y,\qquad M_a=\frac{H_a}{c^2},\qquad
 \Lambda_a=L_\alpha Y^{-1},\qquad R_a=L_\alpha Y.
 \label{v:eq:lectores-longitud-accion}
\end{equation}
El sector nulo se mantiene separado y no se invierte.

\begin{theorem}[Unicidad del coeficiente gravitatorio en la realización radial]
\label{v:thm:gravedad-unica-alpha}
Los lectores anteriores satisfacen
\begin{equation}
 H_a\Lambda_a=\hbar_ac\,I,\qquad
 R_a\Lambda_a=L_\alpha^2I.
 \label{v:eq:dos-relaciones-gravedad}
\end{equation}
Si la lectura física declarada del radio reducido es
$R_a=(G_a/c^2)M_a$, entonces existe un único coeficiente compatible y es
\begin{equation}
 \boxed{G_a=\frac{c^3L_\alpha^2}{\hbar_a}},
 \qquad
 \boxed{\frac{G_a\hbar_a}{c^3}=L_\alpha^2}.
 \label{v:eq:gravedad-unica-alpha}
\end{equation}
\end{theorem}
\begin{proof}
Las definiciones dan
$H_a\Lambda_a=Q_{L,a}L_\alpha I=\hbar_acI$ y
$R_a\Lambda_a=L_\alpha^2I$. Al eliminar $\Lambda_a$,
\[
 R_a=\frac{L_\alpha^2}{\hbar_ac}H_a.
\]
Por otra parte,
$(G_a/c^2)M_a=G_aH_a/c^4$. Como $H_a$ es invertible sobre la fibra
no nula, ambas expresiones coinciden si y sólo si
$G_a=c^3L_\alpha^2/\hbar_a$. La misma cancelación muestra que cualquier
otro $G_a'$ que conserve las dos relaciones satisface $G_a'=G_a$.
\end{proof}

La prueba usa conjuntamente los dos lectores sobre el mismo carácter; no
identifica la compresión de $Y^{-1}$ con la inversa de una compresión. Una
variación de ruta o una eliminación de memoria que actúe de igual modo
sobre $H_a$ y $R_a$ conserva su proporcionalidad y no introduce un nuevo
factor escalar. El análisis dimensional confirma los exponentes
$c^3L_\alpha^2\hbar_a^{-1}$, pero es
\eqref{v:eq:dos-relaciones-gravedad}, y no ese control dimensional, lo que
fija el coeficiente adimensional a uno.

\subsection{La escritura circular como corolario}

El calendario determina $T_{108}=108t_0$ y
$\omega_{108}=2\pi_{\HMT}/T_{108}$. En la realización circular del mismo
lector longitudinal,
\begin{equation}
 \omega_{108}=\frac{c}{L_\alpha},\qquad
 R_{108}=\frac{c}{\omega_{108}}
 =\frac{54}{\pi_{\HMT}}ct_0=L_\alpha,
 \qquad
 t_0=\frac{\pi_{\HMT}L_\alpha}{54c}.
 \label{v:eq:compatibilidad-reloj-longitud}
\end{equation}
La igualdad declara la compatibilidad entre el avance levantado del reloj
y la sección longitudinal; no utiliza $G_a$ para escoger $L_\alpha$.
Sustituirla en \eqref{v:eq:gravedad-unica-alpha} produce
\begin{equation}
 G_a=\left(\frac{54}{\pi_{\HMT}}\right)^2
 \frac{c^5t_0^2}{\hbar_a}.
 \label{v:eq:gravedad-circular-corolario}
\end{equation}
Por consiguiente, en la expresión parametrizada
$G_a(\vartheta)=\vartheta c^5t_0^2/\hbar_a$, el valor
$\vartheta=(54/\pi_{\HMT})^2$ es un corolario de la longitud construida,
no un selector circular independiente.

Con
$E_{108,a}=E_{P,a}=Q_{L,a}=\hbar_a\omega_{108}$ y
$m_{108,a}=E_{108,a}/c^2$ se conservan asimismo
\[
 \frac{\hbar_a}{m_{108,a}c}=R_{108}=L_\alpha,
 \qquad
 \frac{h_a}{m_{108,a}c}=2\pi_{\HMT}L_\alpha.
\]
Finalmente, las dos orientaciones de acción satisfacen
\begin{equation}
 \frac{G_{\mathrm{pre}}}{G_{\mathrm{ret}}}
 =\frac{\hbar_{\mathrm{ret}}}{\hbar_{\mathrm{pre}}}
 =R_{\mathrm{act}}^{-1},
 \qquad
 \ell_P^2:=L_\alpha^2
 =\frac{G_a\hbar_a}{c^3}.
 \label{v:eq:reciprocidad-gravedad-accion}
\end{equation}
La unidad areal procede así de $L_\alpha^2$; la reciprocidad posterior
entre acción y gravedad conserva esa área al cambiar de orientación.
